
lieu de la th\'eorie de Kummer le fait suivant, d\^u \`a Grothendieck (SGA 1 XI, \`a la derni\`ere page)~: si \(X\) est un sch\'ema propre sur un corps \(k\), avec \(H^0(X,{O}_X)=k\), alors il existe un isomorphisme de \(k\)-sch\'emas en groupes \(\underline{\Hom}_{\mathrm{gr}}(\mu_n,\underline{\Pic}_{X/k})\simeq \underline{H}^1(X,\mathbb Z/n\mathbb Z)\); par cons\'equent, si \(k\) est alg\'ebriquement clos, on a \(\Hom_{\mathrm{gr}}(\mu_n,\underline{\Pic}_{X/k})\simeq H^1(X,\mathbb Z/n\mathbb Z)\).
\item Si \(k\) est de caract\'eristique nulle, alors le noyau \(N\) est nul, puisque tout groupe fini unipotent est alors nul (SGA 3 XVII 1.7). On retrouve ainsi \`a peu pr\`es le ``Vanishing theorem'' de Kodaira sous la forme donn\'ee par Mumford [2].
\end{enumerate}

\textbf{Th\'eor\`eme 3.13.} \emph{Soient \(S\) un sch\'ema noeth\'erien et \(X\) un \(S\)-sch\'ema projectif, muni d'un faisceau inversible relativement ample \(H={O}_X(1)\). On suppose que les fibres g\'eom\'etriques de \(X/S\) soient int\`egres et de dimension \(r\). Alors pour une famille \(\mathcal L\) de classes de faisceaux inversibles sur les fibres de \(X/S\), les conditions suivantes sont \'equivalentes}~:
\begin{enumerate}[label=(\roman*)]
\item \emph{\(\mathcal L\) est une famille limit\'ee.}
\item \emph{Dans les polyn\^omes de Hilbert \(\chi(L_K(n))=\sum_{i=0}^{r}a_i\binom{n+i}{i}\) des \(L_K\in\mathcal L\), le coefficient \(a_{r-1}\) reste born\'e et le coefficient \(a_{r-2}\) reste minor\'e.}
\item \emph{Lorsque \(L_K\) parcourt \(\mathcal L\), le degr\'e \(d(L_K)=\langle c_1(L_K)\cdot c_1(H_K)^{r-1}\rangle\) reste born\'e et le nombre d'intersection \(\langle c_1(L_K)^2\cdot c_1(H_K)^{r-2}\rangle\) reste minor\'e.}
\item \emph{Il existe un entier \(N\) tel que l'on ait}
\[
H^{-N}<\mathcal L<H^N
\]
\emph{au sens de la relation d'ordre d\'efinie par la relation d'amplitude\,; c'est-\`a-dire, pour tout \(L_K\in\mathcal L\), \(H_K^{\otimes N}\otimes L_K\) et \(H_K^{\otimes N}\otimes L_K^{-1}\) sont ample.}
\end{enumerate}

